\begin{graphicspathcontext}{{./chapters/ai/imgs/},{./chapters/ai/imgs/auto/},\old}

\sidecite{Chlistalla2011,Lewandowsky2017,Burr2018,Vosoughi2018,Feldman2018,Zuboff2019}
\begin{frame}{{Impacts sociétaux :} démocratie, autonomie, confiance}
	\begin{description}
	\item[Désinformation et manipulation] deepfakes, agents conversationnels générant des fake news à grande échelle
	\vspace{.25cm}
	\item[Surveillance de masse] reconnaissance faciale généralisée, score social
	\vspace{.25cm}
	\item[Érosion de l'autonomie] recommandations addictives (réseaux sociaux), nudges manipulateurs, évaporation des capacités et connaissances acquises
	\vspace{.25cm}
	\item[Fragilisation de la confiance] si on ne peut plus distinguer le vrai du faux, les débats démocratiques s'effondrent
	\vspace{.25cm}
	\item[Systèmes autonomes non coopératifs] course à l'armement algorithmique (trading haute fréquence, drones autonomes)
	\end{description}
	\vspace{.5cm}
	\alertbox{L'IA amplifie ce qu'il y a de meilleur... mais aussi de pire dans nos sociétés}
\end{frame}

\end{graphicspathcontext}

\endinput
